% §3 Design — "The Intelligence Ladder" (S&P 2027 reframe, v4 architecture)
% Draft 2026-10-25. Grounded in docs/ARCHITECTURE.md + verified field deployment
% (docs/FIELD_REVIVAL_2026-10-25.md). All claims trace to code/artifacts listed
% in docs/MASTER_PLAN_2026-10-24.md §Paper.3.

\section{Design}
\label{sec:design}

\subsection{Threat model and design goal}
An operator deploys a deception endpoint that faces opportunistic Internet
adversaries. The adversary controls the commands they type and observes every
byte the endpoint returns; we assume no adversary-side compromise of the
endpoint host itself. The operator's cost is twofold: the \emph{monetary and
latency price} of each AI capability added to the serving loop, and the
\emph{trust price}---every generated artifact is a liability if it leaks real
secrets or contradicts the fiction the operator has committed to.

Our design goal is therefore not ``deceive better with more AI'' but
\emph{price} each increment of intelligence: a capability is wired into the
serving loop only if a measured marginal effect justifies its trust price.
The architecture below makes that pricing possible by construction, because
every candidate response must pass the same deterministic gate regardless of
which AI rung produced it.

\subsection{Verified serving loop}
\label{sec:design:loop}
RAGIN is organized as a \emph{harness--session--sandbox} triad
(\S\ref{sec:impl}). The harness is a stateless orchestration loop; all state
lives in an append-only session log with crash-safe resumption; the sandbox
executes attacker interactions and is the only path on which bytes reach the
attacker. The loop for one attacker command is:

\begin{enumerate}
  \item \textbf{Classify.} A behavioral classifier maps the session history to
  a skill level (novice$\to$apt).
  \item \textbf{Retrieve.} A threat-intelligence retrieval stage (hybrid
  BM25/RRF with a deterministic command-aware retriever) grounds the fiction
  in MITRE-mapped corpus evidence; retrieved scores are clamped and
  provenance-stamped, and attacker-supplied text enters the LLM context only
  inside a labeled untrusted block.
  \item \textbf{Generate.} The deception stage drafts a candidate artifact
  under the session's host fiction (pinned identity, per-session policy).
  \item \textbf{Verify.} A deterministic verifier checks the candidate against
  the fiction graph and a real-shape secret battery
  (\S\ref{sec:design:gate}); wrappers add strict-research and deadline
  enforcement. The verifier is fail-closed: unverifiable candidates are
  discarded, never repaired into service.
  \item \textbf{Approve.} A policy gate approves programmatically at low risk;
  a human decision is required only at the operator-chosen
  \texttt{REQUIRE} level, and an absent approver blocks serving entirely.
  \item \textbf{Commit.} Only committed artifacts are served, and each served
  response carries a serve-path provenance label and matched content hashes.
\end{enumerate}

Two properties matter for the rest of the paper. First, the loop is
\emph{rung-agnostic}: rungs \S\ref{sec:design:ladder} replace the generation
stage but the verifier, the approval gate, and the provenance ledger are
invariant, so cross-rung comparisons measure the rung, not the plumbing.
Second, the loop fails closed: in our adversarial hardening campaign
(\S\ref{sec:eval:hardening}), no probe extracted unverified text on any path.

\subsection{The intelligence ladder}
\label{sec:design:ladder}
On top of the invariant loop we instantiate four rungs of increasing
intelligence, each a drop-in generation stage:
\begin{enumerate}
  \item \textbf{Scripted templates} --- deterministic fiction answers for a
  closed command grammar; zero model cost.
  \item \textbf{Behavioral classifier} --- the rung-1 loop steered by skill
  stratification.
  \item \textbf{Retrieval-augmented planning} --- the rung-2 loop grounded in
  the corpus and evidence gate.
  \item \textbf{Free generation} --- an LLM drafts the artifact under the same
  verification regime.
\end{enumerate}
Decision capability is priced on the same ladder: a deterministic policy
selector is the \emph{decision floor}; learned decision models (a fine-tuned
router, a typed-question judgment model) attach only as bounded, non-authoritative
signals that must clear the same frozen gate battery as any other component.

\subsection{The verification gate}
\label{sec:design:gate}
The gate combines (i) fiction-consistency checks against the committed graph,
(ii) a danger battery that matches \emph{real} credential shapes---GitHub and
Slack tokens, PEM blocks, JWTs, exact-length AWS key IDs---over a
Unicode-normalized view that defeats zero-width and homoglyph smuggling, and
(iii) an allowlist of serve-path templates for deterministic answers. The
battery is adversarially maintained: a two-wave red-team campaign
(\S\ref{sec:eval:hardening}) contributed every pattern, and the deception
fiction's own placeholder syntax is exempt by construction so the honeypot's
bait is never mistaken for a leak.

\subsection{Adversarial workload as measurement}
\label{sec:design:workload}
Passive background traffic cannot answer engagement questions (dwell, artifact
access, credential interaction) at statistical depth. We therefore drive the
live endpoint with a scripted team of LLM-operated attacker personas---novice,
reconnaissance operator, credential hunter, patient APT operator, and an
evasion specialist instructed to detect the honeypot---each issuing
single-command turns over SSH against the production loop. The workload is
reproducible (same persona prompts, per-turn LLM sampling) and comparable
across arms, and it is exactly the measurement the fold-down analysis
requires: what does each additional rung buy \emph{against a fixed adversary}?
